\FloatBarrier
\newpage
\section{Tool to measure the hydrate layer thickness and particle distances}
\label{chap:hydrate-rim}

\textit{Some of the results and data presented in this chapter have already been published in \cite{kleiner.2024a,kleiner.2026b}. However, the datasets were significantly extended and re-evaluated.}

The quantitative analysis of microstructural features from 2D images is fundamental to understand the behaviour of cementitious materials.
Therefore, this section introduces and validates an automated image analysis tool which enables multiple applications.
First, the development of \ac{CSH} on the micro-scale is of interest to understand and later model the cement hydration.
Image-based \ac{HLT} measurements provide a direct, phase-specific assessment of these reaction rates.
The method can also be applied to conduct a particle distance analysis.
This is of interest, since the inter-particle spacing in fresh pastes governs rheological properties.
This analysis directly addresses challenges associated with modern binders, such as the agglomeration tendencies of calcined clays, by objectively quantifying the dispersion efficiency of superplasticisers or sonication.

The methods, which are detailed in \zcref{subsec:algo_hlm} and \zcref{sec:ebsd-hydrate}, were first validated using synthetic datasets.
Subsequently, they were applied to a range of experimental samples, including pure hydrated \ac{C3S} and \ac{C2S}, their mixtures at various ages, and fresh pastes subjected to different dispersion techniques.
Furthermore, the \ac{HLT} measurements were correlated with \ac{EBSD} data to investigate the influence of crystallographic orientation on hydrate growth.
The raw images used for this chapter are publicly available \cite{kleiner.2023a, kleiner.2025e, kleiner.2025d, kleiner.2026a}.
